% !TeX program = pdflatex
%==============================================================================
%  Lista Teórica 04 --- Métodos Não Paramétricos (KNN)
%  O gabarito sai de "Lista teorica 04 - gabarito.tex".
%==============================================================================
\providecommand{\gabaritoopt}{}
\documentclass[11pt,a4paper]{article}
\usepackage[\gabaritoopt]{../../recursos/latex/estilo-lista}

\begin{document}
\cabecalholista{04}{Métodos Não Paramétricos (KNN)}

%------------------------------------------------------------------------------
\begin{exercicio}[o KNN nos dois extremos]
Considere a amostra abaixo, com $n=5$ e uma única covariável:
\begin{center}\small
\begin{tabular}{@{}lrrrrr@{}}
\toprule
$x_i$ & $0$ & $1$ & $3$ & $6$ & $10$ \\
$y_i$ & $2$ & $3$ & $7$ & $5$ & $12$ \\
\bottomrule
\end{tabular}
\end{center}
\begin{enumerate}[label=(\alph*),leftmargin=1.1cm]
  \item Calcule $\rhat(4)$ pelo KNN com $k=1$, $k=3$ e $k=5$.
  \item Para uma amostra genérica sem $x_i$ repetidos, qual é o erro quadrático
        médio de \emph{treino} do KNN com $k=1$? E com $k=n$?
  \item Descreva, em uma frase cada, a cara da função $\rhat$ ajustada nos dois
        extremos.
  \item Qual dos dois extremos tem viés alto e qual tem variância alta? Relacione
        com a Aula~01.
\end{enumerate}
\end{exercicio}

\begin{solucao}
\textbf{(a)} As distâncias a $x_0=4$ são $|0-4|=4$, $|1-4|=3$, $|3-4|=1$,
$|6-4|=2$, $|10-4|=6$. Em ordem: $x=3$ (d.\ $1$), $x=6$ (d.\ $2$), $x=1$ (d.\ $3$),
$x=0$ (d.\ $4$), $x=10$ (d.\ $6$). Logo
\[
  k=1:\ \rhat(4)=7; \qquad
  k=3:\ \rhat(4)=\tfrac{7+5+3}{3}=\mathbf{5}; \qquad
  k=5:\ \rhat(4)=\tfrac{2+3+7+5+12}{5}=\mathbf{5{,}8}.
\]

\smallskip
\textbf{(b)} Com $k=1$ o erro de treino é \textbf{exatamente zero}: o vizinho mais
próximo de $\X_i$ é o próprio $\X_i$ (distância $0$), então
$\rhat(\X_i)=Y_i$ para todo $i$. Com $k=n$ todos os pontos são vizinhos de todos,
$\rhat(\x)\equiv\bar Y$, e o erro de treino é
$\frac1n\sum_i(Y_i-\bar Y)^2$: a \emph{variância amostral} de $Y$ (a menos do
fator $n/(n-1)$).

\smallskip
\textbf{(c)} Com $k=1$, $\rhat$ é uma função \textbf{escada} que passa por todos
os pontos: constante em cada célula de Voronoi, saltando no meio do caminho entre
observações vizinhas. Com $k=n$, $\rhat$ é uma \textbf{reta horizontal} na altura
$\bar Y$.

\smallskip
\textbf{(d)} $k=1$ é o extremo de \textbf{variância alta} e viés baixo: cada
predição depende de uma única observação, e trocar a amostra muda tudo. $k=n$ é o
extremo de \textbf{viés alto} e variância zero: a predição não depende de $\x$, e
por isso não pode acompanhar $r$ --- é literalmente o preditor constante do
Exercício~2 da Lista Teórica~01, cujo risco é $\Var(Y)$. O $k$ é o
\emph{tuning parameter} que percorre o balanço da Aula~01, com a diferença de que
aqui \textbf{mais flexível é $k$ menor}, ao contrário do grau do polinômio.
\end{solucao}

%------------------------------------------------------------------------------
\begin{exercicio}[o KNN é um suavizador linear]
Todos os métodos desta aula se escrevem como $\rhat(\x)=\sum_i w_i(\x)Y_i$ ---
são \textbf{suavizadores lineares}. Avaliando nos próprios pontos de treino,
$\widehat{\bm Y}=\bm H\bm Y$, onde $H_{ij}=w_j(\X_i)$.
\begin{enumerate}[label=(\alph*),leftmargin=1.1cm]
  \item Para o KNN com $k$ vizinhos e $x_i$ distintos, mostre que $h_{ii}=1/k$
        para todo $i$, e conclua que $\operatorname{tr}(\bm H)=n/k$.
  \item A quantidade $\operatorname{tr}(\bm H)$ é o \textbf{número efetivo de
        parâmetros} do ajuste. Calcule-a para $n=100$ e $k=1$, $k=5$, $k=25$.
        Compare com o número de parâmetros de um polinômio de grau $p$ e comente.
  \item Aplique a fórmula do LOOCV da Aula~03,
        $\riscohat_{\text{LOO}}=\frac1n\sum_i\big((Y_i-\rhat(\X_i))/(1-h_{ii})\big)^2$,
        ao KNN. O que acontece quando $k=1$? O resultado faz sentido?
\end{enumerate}
\end{exercicio}

\begin{solucao}
\textbf{(a)} Fixe $i$ e considere a predição no próprio ponto de treino $\X_i$. Os
$k$ vizinhos mais próximos de $\X_i$ \emph{incluem o próprio $\X_i$}, que está a
distância $0$ --- e é o único a essa distância, já que os pontos são distintos. O
KNN dá peso $1/k$ a cada um dos $k$ vizinhos, então o peso da observação $i$ sobre
a própria predição é $h_{ii}=1/k$. Somando na diagonal,
\[
  \operatorname{tr}(\bm H)=\sum_{i=1}^n h_{ii}=n\cdot\frac1k=\frac nk .
\]

\smallskip
\textbf{(b)} Com $n=100$:
\begin{center}\small
\begin{tabular}{@{}lccc@{}}
\toprule
$k$ & $1$ & $5$ & $25$ \\
\midrule
$\operatorname{tr}(\bm H)=n/k$ & $100$ & $20$ & $4$ \\
\bottomrule
\end{tabular}
\end{center}
Com $k=1$ o KNN gasta $100$ parâmetros efetivos para $100$ observações: um
parâmetro por ponto, que é a definição de interpolar --- e casa com o erro de
treino zero do Exercício~1(b). Com $k=25$ ele usa $4$, menos que um polinômio de
grau~4.

A comparação com o polinômio é a moral do exercício: $p+1$ conta parâmetros
\emph{explícitos}, e $\operatorname{tr}(\bm H)$ generaliza isso para métodos que
não têm parâmetros nenhum no sentido usual. É a mesma quantidade nos dois casos
--- para o MQO com $p+1$ colunas, $\operatorname{tr}(\bm H)=p+1$ --- e é ela que
aparece na fórmula do otimismo do Exercício~1 da Lista Teórica~03. Ou seja:
``complexidade'' tem uma definição operacional que não depende de o modelo ser
paramétrico.

\smallskip
\textbf{(c)} Com $k=1$ temos $h_{ii}=1$, logo $1-h_{ii}=0$: a fórmula divide por
zero. E o numerador também é zero, porque $\rhat(\X_i)=Y_i$. Ou seja, $0/0$ --- a
fórmula não diz nada.

E está certa em não dizer. O atalho vale para ajustes lineares em que nenhuma
observação determina sozinha a própria predição; com $k=1$ ela a determina
inteiramente, e o resíduo de treino perdeu \emph{toda} a informação sobre o erro
de validação. Não há como recuperar o LOOCV de um único ajuste nesse caso: é
preciso rodar a força bruta, refazendo o ajuste $n$ vezes. Para $k\ge 2$ a fórmula
funciona normalmente, com o fator de inflação $1/(1-1/k)=k/(k-1)$.
\end{solucao}


\end{document}
